\documentclass[11pt,reqno]{amsart}

\usepackage[T1]{fontenc}
\usepackage{lmodern}
\usepackage{microtype}
\usepackage{mathtools,amssymb,amsfonts}
\usepackage{booktabs}
\usepackage{xcolor}
\usepackage[margin=1.08in]{geometry}
\usepackage[colorlinks=true,linkcolor=blue!55!black,urlcolor=blue!55!black]{hyperref}

\newtheorem{theorem}{Theorem}[section]
\newtheorem{proposition}[theorem]{Proposition}
\newtheorem{lemma}[theorem]{Lemma}
\theoremstyle{remark}
\newtheorem{remark}[theorem]{Remark}

\newcommand{\R}{\mathbb R}
\newcommand{\C}{\mathbb C}
\newcommand{\one}{\mathbf 1}
\newcommand{\SO}{\mathrm{SO}}
\newcommand{\eps}{\varepsilon}
\newcommand{\ip}[2]{\left\langle #1,#2\right\rangle}
\newcommand{\abs}[1]{\left\lvert #1\right\rvert}

\title{A dimension-five reduction for the\\Makeev counterexample}
\author{Working note}
\date{August 2026}

\begin{document}

\begin{abstract}
The dimension-four counterexample suggests using the same function on
$S^4\subset\C^2\oplus\R$,
\[
 f_\eps(z_1,z_2,x)=\operatorname{Re}(z_1^2\overline z_2)
 +\eps\abs{z_1}^2\operatorname{Im}(z_1^2\overline z_2).
\]
This note proves all dimension-independent parts of that extension and reduces
the remaining dimension-five transversality assertion to a unit-ideal
calculation over $\mathbb Q$.  The accompanying Magma file
\texttt{makeev\_d5\_certificate.m} performs that exact calculation.  The
calculation could not be run in the environment in which this note was
prepared, so the theorem below is conditional until the program prints its
stated certificate.  Numerical searches strongly support the assertion and
also support the same candidate in several higher dimensions.
\end{abstract}

\maketitle

\section{Statement and status}

Put
\[
 Q(z_1,z_2,x)=\operatorname{Re}(z_1^2\overline z_2),\qquad
 H(z_1,z_2,x)=\abs{z_1}^2\operatorname{Im}(z_1^2\overline z_2).
\]

\begin{theorem}[Conditional only on the exact certificate]\label{thm:main}
If the saturated ideal constructed in
\texttt{makeev\_d5\_certificate.m} is the unit ideal, then there is
$\eps_0>0$ such that, whenever $0<\abs\eps<\eps_0$, the smooth odd function
$Q+\eps H$ on $S^4$ agrees with no linear functional on any rotated normalized
$A_5$ root system.  Consequently Makeev's conjecture fails in dimension five.
\end{theorem}

The required output is
\begin{verbatim}
Certificate verified: the saturated ideal is the unit ideal.
\end{verbatim}
The ideal has rational coefficients.  Thus this is an exact symbolic
certificate, not a numerical test.

\section{Simplex labels and the cubic annihilator}

Let
\[
 V=\one^\perp\subset\R^6,\qquad
 q_i=e_i-\frac16\one,\qquad
 u_{ij}=\frac{q_i-q_j}{\sqrt2}.
\]
For an odd function $F$ and an isometry $g:V\to\R^5$, put
\[
 Y_F(g)_{ij}=F(gu_{ij}),\qquad Y_F(g)_{ji}=-Y_F(g)_{ij}.
\]
A skew edge-label matrix is called exact if
$Y_{ij}=\rho_i-\rho_j$.  As in the dimension-four argument,
\begin{equation}\label{eq:exact}
 Y\text{ is exact}\quad\Longleftrightarrow\quad
 \Pi Y\Pi=0,\qquad
 \Pi=I-\frac16\one\one^T.
\end{equation}
This is also equivalent to the vanishing of all triangle sums
$Y_{0i}+Y_{ij}+Y_{j0}$, $1\le i<j\le5$.

Identify $\R^5=\C^2\oplus\R$ and set
\[
 J(z_1,z_2,x)=(iz_1,2iz_2,0),\qquad
 h_\theta(z_1,z_2,x)=(e^{i\theta}z_1,e^{2i\theta}z_2,x).
\]
Both $Q$ and $H$ are invariant under $h_\theta$.  Write $p_i=gq_i$ and define
\[
 X(g)_{ij}=\ip{p_i}{Jp_j}_{\R},\qquad
 \Lambda_g(Y)=\sum_{i<j}X(g)_{ij}Y_{ij}.
\]
Since $\sum_i p_i=0$, every row sum of $X(g)$ vanishes.  Hence
$\Lambda_g$ annihilates exact labels.

The cubic identity used in dimension four is independent of the dimension.
If
\[
 P_3(v)=\sum_{i=0}^5\ip{q_i}{v}^3,
\]
then, for a homogeneous cubic $C$ and $K\in\mathfrak{so}(5)$,
\begin{equation}\label{eq:cubicidentity}
 \ip{C}{K\cdot P_3}
 =2\sqrt2\sum_{i<j}\ip{q_i}{Kq_j}_{\R}C(u_{ij}).
\end{equation}
The proof is exactly the Parseval-frame expansion from the dimension-four
paper: $\sum_iq_iq_i^T=I_V$ is the only frame fact used.  Taking
$K=g^{-1}Jg$ and $C=Q\circ g$, the infinitesimal invariance $J\cdot Q=0$
gives
\begin{equation}\label{eq:globalzero}
 \Lambda_g(Y_Q(g))=0\qquad\text{for every }g.
\end{equation}

Thus the only new statement needed in dimension five is
\begin{equation}\label{eq:target}
 Y_Q(g)\text{ exact}\quad\Longrightarrow\quad
 \Lambda_g(Y_H(g))\ne0.
\end{equation}

\section{A nine-equation description of the cubic zeros}

Write
\[
 p_i=(a_i,b_i,c_i)\in\C^2\oplus\R,
 \qquad a,b\in\C^6,\quad c\in\R^6.
\]
The frame identities say that
\[
 \one,a,b,c,\overline a,\overline b
\]
are mutually orthogonal in $\C^6$, with squared norms
$6,2,2,1,2,2$, respectively.  Equivalently,
\begin{gather*}
 \sum a_i=\sum b_i=\sum c_i=0,\qquad
 \sum a_i^2=\sum b_i^2=\sum a_ib_i=\sum a_i\overline b_i=0,\\
 \sum\abs{a_i}^2=\sum\abs{b_i}^2=2,\qquad
 \sum c_i^2=1,\qquad
 \sum c_ia_i=\sum c_ib_i=0.
\end{gather*}

At the edge $gu_{ij}$, the complex cubic has, apart from the common factor
$1/(2\sqrt2)$, the label
\[
 \widetilde Y_{ij}=(a_i-a_j)^2(\overline b_i-\overline b_j).
\]
For $(u\wedge v)_{ij}=u_iv_j-v_iu_j$, expansion modulo exact matrices gives
\begin{equation}\label{eq:omega}
 \omega:=\Pi\widetilde Y\Pi
 =-a^2\wedge\overline b+2a\wedge(a\overline b).
\end{equation}
The $Q$-labels are exact precisely when
$\operatorname{Re}\omega=0$, or $\omega+\overline\omega=0$.

Introduce
\begin{align*}
 s&=\frac12\sum\abs{a_i}^2a_i,&
 m&=\sum a_i^2\overline b_i,&
 r&=\sum c_i a_i^2,\\
 n&=\sum a_i^3,&
 \ell&=\frac12\sum a_i^2b_i,&
 \lambda&=\frac12\sum\abs{a_i}^2\overline b_i,\\
 u&=\frac12\sum a_i\overline b_i^2,&
 w&=\sum c_i a_i\overline b_i,&
 d&=\frac12\sum a_i\abs{b_i}^2.
\end{align*}
Expansion in the displayed orthogonal basis gives
\begin{align*}
 a^2&=sa+\frac m2b+rc+\frac n2\overline a+\ell\overline b,\\
 a\overline b&=\lambda a+ub+wc+\frac m2\overline a+d\overline b.
\end{align*}
Substitution in \eqref{eq:omega} yields
\begin{align*}
 \omega={}&2u\,a\wedge b+2w\,a\wedge c+m\,a\wedge\overline a
 +(2d-s)a\wedge\overline b\\
 &-\frac m2b\wedge\overline b-r\,c\wedge\overline b
 -\frac n2\overline a\wedge\overline b.
\end{align*}
The wedges of the five vectors $a,b,c,\overline a,\overline b$ are linearly
independent.  Comparing the last display with its conjugate proves the
following exact criterion.

\begin{proposition}[Cubic exactness criterion]\label{prop:criterion}
The $Q$-labels are exact if and only if
\begin{equation}\label{eq:nine}
 \operatorname{Im}m=0,\qquad
 4u=\overline n,\qquad
 r=0,\qquad w=0,\qquad 2d=s.
\end{equation}
These are nine real cubic equations.
\end{proposition}

This criterion is the key simplification: there is no need to use 25 entries
of an orthogonal matrix, nor to use a rational parametrization of
$\SO(5)$.

There is also a useful dimension-independent version.  For any number $N$ of
simplex vertices, exactness forces the components of $a^2$ and
$a\overline b$ perpendicular to
$\operatorname{span}_{\C}\{\one,a,b,\overline a,\overline b\}$ to vanish.
Together with the five coefficient equations in
\eqref{eq:nine}, this is equivalent to the two coordinatewise vector
identities
\begin{align}
 a^2&=sa+\frac m2b+\frac n2\overline a+\ell\overline b,
 \label{eq:universal-a2}\\
 a\overline b&=\lambda a+\frac{\overline n}{4}b
 +\frac m2\overline a+\frac s2\overline b.
 \label{eq:universal-ab}
\end{align}
Here $m$ is real.  Conversely, the two identities imply exactness by
substitution in \eqref{eq:omega}.  Equations
\eqref{eq:universal-a2}--\eqref{eq:universal-ab} therefore give a
quadratic, linearly growing system for testing every fixed $d\ge4$.  The file
\texttt{makeev\_general\_certificate.m} implements that system.

\section{The exact finite calculation}

For completeness, define the rational polynomial whose nonvanishing is needed.
For $d_{ij}=a_i-a_j$ and $e_{ij}=b_i-b_j$, put
\begin{align}
 \mathcal L={}&\sum_{i<j}
 \bigl(\operatorname{Im}(a_i\overline a_j)
      +2\operatorname{Im}(b_i\overline b_j)\bigr)
 \abs{d_{ij}}^2
 \operatorname{Im}(d_{ij}^2\overline e_{ij}).
 \label{eq:L}
\end{align}
All real and imaginary parts in \eqref{eq:L} are polynomials in the real
coordinates.  Directly from the normalization of the roots,
\begin{equation}\label{eq:Lscale}
 \Lambda_g(Y_H(g))=\frac{\mathcal L}{4\sqrt2}.
\end{equation}

The Magma program makes the following exact choices.
\begin{enumerate}
\item It writes $a_i=x_i+iy_i$, $b_i=s_i+it_i$ and eliminates the last
coordinate in each of the five zero-sum real vectors $x,y,s,t,c$.
\item The circle symmetry sends $(a,b,c)$ to
$(e^{i\theta}a,e^{2i\theta}b,c)$.  Since $a\ne0$, a vertex permutation and
this symmetry put every orbit in the chart $y_0=0$, $x_0\ne0$.
\item It imposes the ten real frame equations, the five equations making $c$
the unit orthogonal complement, and the nine equations
\eqref{eq:nine}.  After the gauge choice these are 24 equations in 24
variables.
\item An inverse variable with equation $z x_0-1=0$ saturates by the chart
condition $x_0\ne0$.  Finally it adds $\mathcal L=0$.
\end{enumerate}
The program asks whether the resulting ideal over $\mathbb Q$ is proper.  If
the answer is no, there is no complex solution, hence in particular no real
configuration satisfying both cubic exactness and $\mathcal L=0$.  This proves
\eqref{eq:target}.  The gauge loses no real orbit: some $a_i$ is nonzero,
vertex permutations preserve all the equations, and the circle action makes
that coordinate real and nonzero.

\section{Completion after the certificate}

Assume the unit-ideal output.  If arbitrarily small nonzero $\eps$ admitted
exact labels for $Q+\eps H$, take $\eps_\nu\to0$ and corresponding rotations
$g_\nu$.  Compactness gives a subsequence $g_\nu\to g$.  Passing to the limit
in the exact-label defect shows that $Y_Q(g)$ is exact.  Since $\Lambda_{g_\nu}$
annihilates exact labels and \eqref{eq:globalzero} holds identically,
\[
 \Lambda_{g_\nu}(Y_H(g_\nu))=0.
\]
The limit contradicts \eqref{eq:target}.  This proves
Theorem~\ref{thm:main}.

An orientation issue does not arise: a transposition of two simplex vertices
induces an orientation-reversing symmetry of the $A_5$ root system.  Therefore
the $\mathrm O(5)$ and $\SO(5)$ formulations are equivalent.

\section{Evidence in higher dimensions}

The same formula makes sense on
$S^{d-1}\subset\C^2\oplus\R^{d-4}$ for every $d\ge4$, and the entire argument
through \eqref{eq:globalzero} remains valid.  Only the classification of the
cubic exact-label locus changes.

There is a distinguished cubic zero in every dimension.  With
$N=d+1$, $\zeta=e^{2\pi i/N}$, take
\[
 a_j=\sqrt{\frac2N}\,\zeta^j,\qquad
 b_j=\sqrt{\frac2N}\,\zeta^{2j}.
\]
The remaining $d-4$ real frame vectors can be completed orthonormally.  Since
\[
 (1-\zeta^r)^2(1-\overline\zeta^{,2r})
\]
is purely imaginary, every $Q$-edge label is zero.  At this zero one obtains,
for $N\ge6$,
\[
 \Lambda(Y_H)=-\frac{3}{N^{3/2}}\ne0.
\]
This checks one orbit, not all of them.

The accompanying Python program searches numerically for other cubic zeros.
The following reconnaissance used independent random starts; failed starts are
excluded from the minimum.
\begin{center}
\begin{tabular}{c@{\qquad}c@{\qquad}c}
\toprule
$d$&converged starts&smallest sampled $\abs{\Lambda(Y_H)}$\\
\midrule
5&99/100&$0.037037037037\ (=1/27)$\\
6&10/10&$0.033452928407$\\
7&4/5&$0.012865396105$\\
8&8/10&$0.008659274357$\\
9&9/10&$0.009496590853$\\
10&6/10&$0.019556538182$\\
\bottomrule
\end{tabular}
\end{center}
No sampled zero had vanishing obstruction.  This is useful evidence that the
same perturbation may work beyond dimension five, but it is not a proof and
does not supply a uniform bound as $d$ varies.

For an exact test in another fixed dimension, change the first assignment
\texttt{d := 5;} in \texttt{makeev\_general\_certificate.m}.  The program
uses \eqref{eq:universal-a2}--\eqref{eq:universal-ab} and a cubic expression
for the obstruction.  A unit-ideal output proves nonvanishing in that selected
dimension; nontermination is merely an inconclusive computation.

\end{document}
